\subsection{Milestone 1 -- Costruzione del dataset}

\subsubsection{Versioni e ticket Jira}

Le versioni scaricate tramite Jira REST API sono state filtrate per: campo
\texttt{released = true} e presenza di \texttt{releaseDate}. Sono emerse 42 versioni
valide per Apache OpenJPA. Per la costruzione del dataset è stato considerato solo il
\textbf{primo 33\% delle release} in ordine cronologico (14 su 42), per ridurre al minimo
i bug dormienti (\textit{snoring bugs}) che compaiono statisticamente nelle versioni più
recenti. Le release saltate (3 versioni senza tag Git) non compensano il calcolo del 33\%,
per la stessa ragione.

I ticket Jira sono stati scaricati con il filtro JQL
\textit{``issuetype = Bug AND status in (Closed, Resolved) AND resolution = Fixed''} e
modellati con i campi: \texttt{created}, \texttt{affectedVersions} (AV),
\texttt{fixVersions} (FV), \texttt{openingVersion} (OV = prima versione precedente alla
creazione), \texttt{injectedVersion} (IV = AV più antica).

\subsubsection{Arricchimento e check di consistenza}

Prima dei check, ogni ticket subisce: (1) calcolo della OV; (2) riduzione delle FV a una
sola (la più recente, ritenuta più affidabile); (3) rimozione delle versioni senza
\texttt{releaseDate}; (4) se AV assenti ma FV disponibili, le FV tranne l'ultima vengono
interpretate come AV --- approccio preferito a Proportion perché sfrutta l'informazione
già presente nel ticket, pur restando un'ipotesi. I ticket privi di data di creazione o
antecedenti a qualsiasi versione vengono scartati (qualità dati preferita alla quantità).

I check di consistenza verificano, per ogni ticket: presenza di chiave, data di creazione,
almeno una FV e una AV con data di rilascio, FV temporalmente successiva a ogni AV, OV
precedente o uguale alla FV. Per le versioni: presenza di nome, data di rilascio e stato
\textit{released}. I ticket non conformi vengono rimossi; le versioni inconsistenti vengono
eliminate anche dai riferimenti nei ticket.

\subsubsection{Stima AV con Proportion}

Per i ticket ancora privi di AV dopo l'arricchimento, le AV vengono stimate tramite
\textbf{Proportion} con approccio \emph{cold start}:
\[
P = \text{mean}\!\left(\frac{FV - IV}{FV - OV}\right), \qquad IV_{\text{stimata}} = FV - P \cdot (FV - OV)
\]
I ticket con $FV = OV$ o con IV stimata $= FV$ vengono scartati.

\subsubsection{Estrazione snapshot Git}

Per ogni release viene identificato il tag Git con match esatto (es.\ \texttt{1.0.0}) o
con suffisso dopo trattino (es.\ \texttt{0.9.6-incubating}), escludendo tag SNAPSHOT. Tre
versioni Jira (\texttt{0.9.0}, \texttt{2.0.0-M1}, \texttt{2.0.0-M2}) non hanno tag Git
(il progetto era su SVN in quel periodo) e vengono saltate, mantenendo però i relativi ticket.

I file vengono elencati con \texttt{git ls-tree -r --name-only <tag>} e il sorgente estratto
con \texttt{git show <tag>:<path>}. I file inclusi devono: terminare con \texttt{.java},
non contenere directory di test (\texttt{/test/}, \texttt{/tests/}), non terminare con
\texttt{Test.java} o \texttt{Tests.java}.

\paragraph{Scelta: full clone vs shallow clone.}
La prima versione della pipeline eseguiva uno \emph{shallow clone} (\texttt{depth=1}),
dimezzando i tempi di download ma causando un bug silenzioso: \texttt{CommitIndexService}
e \texttt{GitLogStatsService}, che scansionano l'intera storia dei commit, vedevano meno
del 10\% dei commit reali. Il tasso di bugginess risultante era $\approx$60\% invece del
$\approx$9\% atteso. La correzione è il \textbf{full clone} via JGit, che garantisce
l'intera storia disponibile.

\subsubsection{Bugginess labeling}

Uno snapshot (\emph{classPath}, \emph{release}) è etichettato \textbf{buggy} se esiste
almeno un ticket tale che:
\begin{enumerate}
    \item la release è tra le AV del ticket;
    \item almeno un commit che cita la chiave del ticket tocca quel \emph{classPath}.
\end{enumerate}
Il linkage commit--ticket è costruito tramite un singolo
\texttt{git log --all --name-only}, che associa a ogni chiave di ticket i file toccati
dai commit che la citano.

\subsubsection{Metriche}

Vengono calcolate 19 metriche per ogni snapshot. Le \textbf{metriche puntuali} sono
calcolate sullo stato del file all'ultimo commit della release: \textbf{LOC} (righe non
vuote) e \textbf{Previous\_Release\_Code\_Smells}. Le \textbf{metriche cumulative} sono
calcolate su tutti i commit del file dall'inizio della storia fino alla release:

\begin{itemize}[noitemsep]
  \item \textbf{LOC\_Touched}: $\sum(\text{righeAgg} + \text{righeCan})$ --- volume totale di riscrittura.
  \item \textbf{NR}: numero totale di commit sul file.
  \item \textbf{Nfix}: commit che citano un ticket con $FV = \text{release corrente}$ e toccano il file.
  \item \textbf{Nauth}: autori distinti (per email).
  \item \textbf{LOC\_Added/Max/Avg}: somma, massimo e media delle righe aggiunte per commit.
  \item \textbf{Churn/Max/Avg}: righe aggiunte + cancellate per commit (riscrittura netta).
  \item \textbf{Change\_Set\_Size/Max/Avg}: file committati insieme (accoppiamento strutturale).
  \item \textbf{Age}: giorni dal primo commit alla data di rilascio.
  \item \textbf{Weighted\_Age}: $\text{Age} \times \text{LOC\_Touched}$ (anzianità ponderata per attività).
  \item \textbf{EXP} e \textbf{SEXP}: esperienza autore in commit-day distinti, globale vs ristretta al sotto-sistema. Calcolate via ricerca binaria su indici preordinati costruiti \textit{una tantum}.
\end{itemize}

\paragraph{Code smell e shift temporale.}
I code smell sono rilevati tramite \textbf{PMD} con i ruleset \texttt{bestpractices},
\texttt{design}, \texttt{errorprone} e \texttt{codestyle}. La metrica è denominata
\textbf{Previous\_Release\_Code\_Smells} perché i smell introdotti nella release $N$ non
impattano $N$ ma le successive:
\[
\text{Previous\_Release\_Code\_Smells}(\text{classe},\, N) = \text{CodeSmells}(\text{classe},\, N-1)
\]
Snapshot della prima release e classi nuove ricevono valore 0. Nel caso di branch paralleli
(branch \texttt{1.x} e \texttt{2.x} di OpenJPA), le classi assenti nella release precedente
ricevono 0: comportamento corretto e atteso rispetto alla logica implementata.

\subsubsection{Aggiornamenti durante la Milestone 2}

Clone automatico via JGit (\texttt{repocloner} package, con fix shallow clone); parametri
PMD estratti da costanti a \texttt{config.yaml} (Tabella~\ref{tab:m1-tuning}); cartelle
output rinominate con prefissi numerici (Tabella~\ref{tab:m1-folders}).
